% ============================================================================
%  OmniWatch Enterprise AI-SOC — Architecture & Performance Validation Report
%  Version 2 · May 2026
%  Compile with: pdflatex main.tex (×2 for TOC)
% ============================================================================
\documentclass[12pt,a4paper]{article}

% ── Packages ────────────────────────────────────────────────────────────────
\usepackage[margin=1in]{geometry}
\usepackage{graphicx}
\usepackage{amsmath,amssymb}
\usepackage{booktabs}
\usepackage{xcolor}
\usepackage{listings}
\usepackage{fancyhdr}
\usepackage{titlesec}
\usepackage{parskip}
\usepackage[hidelinks]{hyperref}

% ── Code listing style ──────────────────────────────────────────────────────
\definecolor{codebg}{HTML}{F5F5F5}
\definecolor{codeframe}{HTML}{CCCCCC}
\lstset{
  basicstyle=\ttfamily\small,
  backgroundcolor=\color{codebg},
  frame=single,
  rulecolor=\color{codeframe},
  breaklines=true,
  columns=fullflexible,
  keepspaces=true,
  xleftmargin=6pt,
  xrightmargin=6pt,
  aboveskip=8pt,
  belowskip=8pt,
}

% ── Header / Footer ────────────────────────────────────────────────────────
\pagestyle{fancy}
\fancyhf{}
\fancyhead[L]{\small OmniWatch AI-SOC --- Validation Report v2}
\fancyhead[R]{\small UMPSA --- Faculty of Computing}
\fancyfoot[C]{\thepage}
\renewcommand{\headrulewidth}{0.4pt}

% ── Section formatting ─────────────────────────────────────────────────────
\titleformat{\section}{\Large\bfseries}{\thesection.}{0.5em}{}
\titleformat{\subsection}{\large\bfseries}{\thesubsection}{0.5em}{}

% ============================================================================
\begin{document}

% ── Title Page ──────────────────────────────────────────────────────────────
\begin{titlepage}
\centering
\vspace*{1.5cm}

{\large\textbf{Universiti Malaysia Pahang Al-Sultan Abdullah}}\\[4pt]
{\large Faculty of Computing}\\[2cm]

\rule{\textwidth}{1pt}\\[0.6cm]
{\LARGE\bfseries OmniWatch Enterprise AI-SOC:\\[6pt]
Architecture \& Performance\\[6pt]
Validation Report (v2)}\\[0.4cm]
\rule{\textwidth}{1pt}\\[1cm]

{\large Internal Technical Documentation}\\[0.4cm]
{\large Final Year Project Report}\\[2cm]

\begin{tabular}{ll}
\textbf{Supervised by:} & Puan Wan Nurulsafawati Binti Wan Manan \\[12pt]
\textbf{Prepared by:}   & SAEED MOHAMMED KHALED ABDULAZIZ -- CF23006 \\
                         & BA NAFEA ABDULLAH HADDAD OBAID -- CB23011 \\
                         & ALHAMED MOHAMMED ABDULLAH MOHAMMED -- CB23010 \\
\end{tabular}

\vfill
{\large May 2026}

\end{titlepage}

% ── Abstract ────────────────────────────────────────────────────────────────
\section*{Abstract}
\addcontentsline{toc}{section}{Abstract}

OmniWatch is an Enterprise AI Security Operations Centre (AI-SOC) platform
engineered for governed post-detection response, auditability, and analyst
supervision in Operational Technology and Industrial Control Systems (OT/ICS)
environments. The platform implements a novel \emph{Split-Brain} architecture:
a stateful Python host performs floating-point machine-learning baseline
computation, while a stateless Rust guest program executes inside a
RISC Zero zero-knowledge virtual machine (zkVM), performing purely
integer threat evaluation on the FPU-less RV32IM instruction set
architecture and producing a cryptographically unforgeable STARK proof
of that computation. This approach ensures verifiable OT/ICS threat detection
and post-detection governance, preventing automated remediation actions from
triggering without mathematical proof of honesty.

\newpage
\tableofcontents
\newpage

% ════════════════════════════════════════════════════════════════════════════
\section{Introduction}
% ════════════════════════════════════════════════════════════════════════════

OmniWatch addresses a critical gap in industrial cyber-security: the
absence of verifiable, auditable post-detection governance for OT/ICS
environments. Modern Security Operations Centres rely on automated SOAR
(Security Orchestration, Automation and Response) workflows that can
execute remediation actions---firewall blocks, service restarts,
isolation commands---without cryptographic proof that the triggering
detection was computed honestly over authentic telemetry.

The platform serves three simultaneous mandates:

\begin{enumerate}
  \item \textbf{Academic benchmark} --- quantifying the latency overhead
        of zero-knowledge virtual machines at the network edge.
  \item \textbf{Enterprise AI-SOC} --- an end-to-end pipeline ingesting
        large-scale real-world datasets (CIC-IDS-2017, BOTSv3) through
        three ML detection pillars with session-scoped SQLite storage,
        SHA-256 chain integrity, and a React dashboard.
  \item \textbf{Cyber-physical testbed} --- a hardware diorama connecting
        an ESP32 Modbus TCP server, Raspberry Pi attack and sensor nodes,
        and a command-node workstation to form a verifiable OT threat-detection loop.
\end{enumerate}

This research provides a substantial contribution to \textbf{SDG~9} (Industry,
Innovation and Infrastructure) through its novel zero-knowledge verification
pipeline for industrial control systems, and \textbf{SDG~11} (Sustainable Cities
and Communities) through its applicability to critical infrastructure governance.

% ════════════════════════════════════════════════════════════════════════════
\section{Bounded Threat Model and Limitations}
% ════════════════════════════════════════════════════════════════════════════

The OmniWatch architecture incorporates several strict operational
limitations that define how the entire system must be interpreted.
These are not merely ``due diligence'' details, but first-class design
constraints that shape the platform's deployment posture.

\subsection{Computational vs. Sensing Integrity}

The trust boundary explicitly terminates at the Raspberry Pi~4 edge node.
The system lacks TPM~2.0 remote attestation and Secure Boot; consequently,
the Pi~4 operating system is the primary attack vector for pre-proof data
manipulation. The architecture achieves strict \textbf{Computational
Integrity} over supplied data (proving the detection algorithms were
executed honestly), but does not guarantee \textbf{Sensing Integrity}
(proving the data genuinely represents the physical reality of the OT network).

\subsection{Static Guest Recompilation Trade-off}

Deterministic threat rules are hardcoded directly into the Rust Guest ELF.
Adding or modifying rules requires a full recompile and reprovisioning
of the zkVM guest image ID (\texttt{VERIFIER\_GUEST\_ID}). This is an accepted
architectural trade-off, preferring absolute state immutability over the
flexibility of hot-reloading rules.

\subsection{Pedagogical Physical Simulation}

The cyber-physical testbed utilizes an ESP32 operating Modbus TCP over Wi-Fi.
However, NCA OTCC-1:2022 mandates strict restriction of wireless technologies
in production OT/ICS networks. This physical setup is therefore classified as
a pedagogical simulation only, demonstrating the verification loop without
meeting the physical layer fidelity required for true SCADA deployments.

To mitigate 2.4\,GHz Wi-Fi spectrum saturation commonly experienced at exhibition venues, the pipeline includes a \textbf{Deterministic Replay Fallback} mode. This accepts pre-recorded PCAP captures and allows out-of-band truncation of the SQLite Spent-Receipt Registry. This is an accepted procedural bypass explicitly scoped to pedagogical demonstration contexts.

% ════════════════════════════════════════════════════════════════════════════
\section{System Architecture}
% ════════════════════════════════════════════════════════════════════════════

\subsection{Hardware Partitioning and Physical Simulation}

The infrastructure uses strict physical service isolation to prevent
computational bottlenecks across the cyber-physical loop.

\begin{table}[htbp]
\centering
\caption{Cyber-Physical Testbed Hardware Allocation}
\label{tab:hardware}
\begin{tabular}{@{}llp{7cm}@{}}
\toprule
\textbf{Node} & \textbf{Hardware} & \textbf{Role} \\
\midrule
OT Target         & ESP32 Type-C Wi-Fi/BT & Modbus TCP server simulating SCADA;
                                             drives 5\,V pumps and LED indicators \\
Attacker           & Pi Zero 2W + TP-Link TL-WN722N & Executes FC\,05 (coil write) and
                                                       FC\,16 (register block write) payloads \\
Network Sensor     & Pi 4 Model B (128\,GB) & Passive TAP listener running Zeek +
                                              CISA ICSNPP Modbus; serialises flows
                                              to bincode \\
Command Node       & Local workstation & RISC Zero STARK prover, React/WASM
                                         verification, FastAPI backend, Ollama LLM \\
\bottomrule
\end{tabular}
\end{table}

\subsection{The Split-Brain Execution Model}

The central architectural invariant is a hard execution boundary between
the Python host and the Rust zkVM guest:

\begin{itemize}
  \item \textbf{Python Host} --- Stateful; accumulates baselines across
        chunks using unrestricted floating-point arithmetic
        (\texttt{numpy}, \texttt{sklearn}, DDSketch). Outputs ML
        parameters: \texttt{p99}, \texttt{mean}, \texttt{stddev},
        and trained LODA/CUSUM models.
  \item \textbf{Rust Guest (zkVM)} --- Stateless; evaluated independently
        per record on the RV32IM ISA with no hardware FPU. Commits a
        \texttt{ThreatVerdict} to the STARK journal containing
        \texttt{input\_hash}, \texttt{is\_threat}, \texttt{category},
        and \texttt{confidence\_pct}.
\end{itemize}

\subsection{Q14 Fixed-Point Encoding}

To bridge the execution boundary, all floating-point ML outputs are
translated into a strict 14-bit fractional integer representation:

\begin{equation}
  Z = \lfloor x \times 2^{14} \rfloor = \lfloor x \times 16384 \rfloor
\end{equation}

Addition and subtraction require no correction. Multiplication requires
a right-shift with rounding:

\begin{equation}
  Z_{\text{product}} = \frac{Z_a \times Z_b + 2^{13}}{2^{14}}
\end{equation}

All Q14 multiplications use Rust's \texttt{saturating\_mul} to prevent
integer overflow from producing incorrect threat verdicts inside the zkVM.

% ════════════════════════════════════════════════════════════════════════════
\section{Data Ingestion and Serialization}
% ════════════════════════════════════════════════════════════════════════════

\subsection{Polars Streaming Pipeline and Bounded Memory}

Large CSV datasets are processed through a \textbf{Polars batched CSV reader}
that eliminates out-of-memory risk. The implementation uses
\texttt{pl.read\_csv\_batched()} with a batch size of 10\,000 rows.
Only one batch resides in RAM at a time, maintaining a peak footprint of
approximately 50--200\,MB regardless of file size. A pandas chunked reader
serves as a fallback when Polars is unavailable.

\subsection{Edge Bincode Serialization}

The Pi~4 edge node serialises each captured Modbus flow into a
\texttt{TelemetryInput} struct encoded in bincode (little-endian, no
padding). This eliminates JSON parsing overhead inside the zkVM, reducing
proof-generation cycle counts.

% ════════════════════════════════════════════════════════════════════════════
\section{Detection Pipeline}
% ════════════════════════════════════════════════════════════════════════════

The detection engine comprises three pillars, specifically structured to
minimise multiplication instructions on the FPU-less RV32IM ISA.

\subsection{DDSketch Volumetric Baselining}

The Python host's \texttt{DDSketchBaseliner} accumulates per-flow
bytes/s observations across all pipeline chunks, returning a Q14-encoded
99th-percentile threshold: $T = \lfloor p_{99} \times 2^{14} \rfloor$.
The Rust guest compares \texttt{current\_scaled = bytes\_per\_sec}
$\times 16384$ against $T$, requiring no floating-point arithmetic.

\subsection{LODA Multivariate Anomaly Detection}

LODA replaces Isolation Forest to bypass cache-hostile pointer jumps and
multiplication-heavy tree nodes. LODA uses sparse projection vectors
with weights $w_{ij} \in \{-1, 0, 1\}$, compiling down to single-cycle
\texttt{ADD/SUB} instructions:

\begin{lstlisting}[language=Rust,caption={LODA zero-MUL projection}]
match model.projections[i * nf + j] {
    1  => z += features[j],   // ADD
    -1 => z -= features[j],   // SUB
    _  => {}
}
\end{lstlisting}

\subsection{Zero-MUL CUSUM Beaconing Detection}

The CUSUM pipeline detects periodic C2 beaconing using only \texttt{ADD},
\texttt{SUB}, and bitwise operations. The DAS-CUSUM accumulator proves:

\begin{equation}
  S_n[i] = \max\!\big(0,\; S_{n-1}[i] + B_n[i] - k\big)
\end{equation}

where $B_n[i]$ is the Q14-encoded L1 score and $k$ is the drift parameter.
The guest instruction count per period relies on \texttt{saturating\_add},
\texttt{saturating\_sub}, and \texttt{.max(0)}---achieving zero multiplications.

\subsection{zkVM Threat Rule Evaluation Matrix}

In addition to the ML pillars, the Rust guest evaluates a deterministic set of threat rules against the \texttt{TelemetryInput}. The evaluation yields a 14-bit flag matrix committed to the STARK journal. Key flags include:

\begin{table}[htbp]
\centering
\caption{Selected zkVM Threat Evaluation Rules}
\label{tab:threat_rules}
\begin{tabular}{@{}llp{8cm}@{}}
\toprule
\textbf{Bit} & \textbf{Flag} & \textbf{Trigger Condition} \\
\midrule
0  & \texttt{HIGH\_RATE} & Absolute bytes/s $> 1$\,MB (no baseline required) \\
2  & \texttt{BRUTE\_FORCE} & Port $\in \{21, 22\}$ $\wedge$ packet count $>$ threshold \\
5  & \texttt{EXFIL\_PATTERN} & Outbound $\wedge$ bytes/s $> 500$\,KB $\wedge$ duration $> 5$\,s \\
8  & \texttt{MODBUS\_WRITE} & Modbus function code $\in \{5, 6, 15, 16\}$ \\
11 & \texttt{DDSKETCH\_VOLUME} & Q14-scaled bytes/s $>$ \texttt{ddsketch\_threshold\_fp14} \\
12 & \texttt{LODA\_ANOMALY} & Sparse projection sum $>$ \texttt{anomaly\_threshold\_fp10} \\
13 & \texttt{CUSUM\_BEACON} & DAS-CUSUM $S_n >$ threshold on any candidate period \\
\bottomrule
\end{tabular}
\end{table}

% ════════════════════════════════════════════════════════════════════════════
\section{Cryptographic Verification and Governance}
% ════════════════════════════════════════════════════════════════════════════

\subsection{STARK Proof Flow}

The system generates raw STARK receipts to prove computational integrity.
The project accepts larger proof sizes ($\sim$217--250\,KB) rather than using
optional Groth16 SNARK wrappers. While SNARK compression would drastically
reduce proof size, it introduces a one-time trusted-setup ceremony. Rejecting
this preserves a completely \textbf{trusted-setup-free} cryptographic environment,
a deliberate trade-off suitable for local high-speed network deployments.

\subsection{Dual-Verification SOAR Gate}

SOAR remediation requires a dual-factor cryptographic gate:

\begin{enumerate}
  \item A machine-generated STARK receipt ($\sim$217--250\,KB).
  \item A human-generated FIDO2/WebAuthn assertion (ECDSA P-256).
\end{enumerate}

To satisfy WebAuthn's strict Relying Party ID requirements in a testbed environment, the command node utilizes a local DNS mapping (\texttt{omniwatch.local} $\rightarrow$ \texttt{127.0.0.1}) secured by locally trusted \texttt{mkcert} TLS certificates. Only if both validations pass concurrently is remediation authorised, ensuring cryptographically signed human-in-the-loop governance.

\subsection{SHA-256 Chain and Spent-Receipt Race Prevention}

The Spent-Receipt Registry uses SQLite's \texttt{INSERT OR IGNORE}
pattern with WAL journal mode to prevent Time-of-Check-to-Time-of-Use
(TOCTOU) replay attacks. This forms the base of a Triple-Layer Audit Trail,
supported by a SHA-256 hash chain ensuring continuous database row integrity,
and STARK proofs verifying the evaluation logic.

% ════════════════════════════════════════════════════════════════════════════
\section{Frontend and Session Management Architecture}
% ════════════════════════════════════════════════════════════════════════════

The frontend architecture is built on React, utilizing TanStack Query for
state management and React Three Fiber for WebGL visualisations.

\subsection{Session-Scoped State and Data Switching}

Active session identity is held in the application's root state and propagated
throughout the component tree. TanStack Query caching uses the \texttt{session\_id}
as part of its query key array. When a new Polars ingestion pipeline completes,
the session ID updates, triggering an automatic cache invalidation and UI
refresh without requiring a hard page reload.

The \texttt{LogExplorer} component implements dynamic data source switching,
allowing the same UI to render live pipeline alerts (via the \texttt{telemetry\_alerts}
table) or legacy static datasets (CIC-IDS-2017, BOTSv3) by mapping different
backend schemas into a canonical 5-tuple \texttt{CicidsLog} shape.

\subsection{Splunk-Tier Network Telemetry View}

The primary analyst interface presents network flows as a dense 5-tuple
(Source IP/Port, Destination IP/Port, Protocol). Following Ecological Interface
Design (EID) principles, row interactions provide progressive disclosure: clicking
a row expands a detail panel containing the raw JSON feature snapshot and
an ``Analyze with AI'' action, which branches into the LLM incident report workflow.

\subsection{Asynchronous WASM Verification}

The WebAssembly module (\texttt{risc0-zkvm}) is instantiated in a dedicated
background Web Worker on the client browser. This serves as an EID UX
optimization: client-side verification updates the visual state asynchronously
without blocking the main thread. However, the authoritative verification is
always strictly enforced via a server-side Python-to-Rust invocation.

\subsection{Spatial WebGL Trust-Chain DAG}

The verification state is visualized in React Three Fiber as a 3D Directed
Acyclic Graph (DAG) mapping the cryptographic trust chain from Raw Telemetry
to STARK Receipt and FIDO2 ECDSA. Each node transitions from a pending state
to a verified state, using shape and color redundancy to maintain EID
accessibility under deuteranopia simulation. If the WebGL context fails to
initialize, the interface degrades gracefully to a 2D CSS-grid matrix.

% ════════════════════════════════════════════════════════════════════════════
\section{LLM RAG Contextualization}
% ════════════════════════════════════════════════════════════════════════════

The platform integrates a local Large Language Model (Phi-3-Mini) to serve as a
translation and contextualization engine for security analysts. Crucially,
the LLM has zero normative weight in the threat detection verdict, which is
exclusively determined by the mathematical zkVM evaluation.

\subsection{Embedding Precision Constraints}

While the generation model operates using INT4 quantization for memory efficiency
on CPU deployments, the embedding model is strictly locked to FP32 or FP16.
Mixed-precision retrieval (where embeddings are aggressively quantized)
degrades cosine similarity accuracy due to quantization noise accumulating
in high-dimensional dot products. This constraint ensures retrieval precision
remains above the calibrated threshold.

\subsection{Source Grounding Enforcement}

To mitigate LLM hallucination risks in critical OT contexts, the UI enforces
strict visible source grounding. The exact document excerpt or raw JSON telemetry
snapshot is displayed alongside the LLM summary. If retrieval confidence falls
below the operational threshold, the UI enters a hard ``No Grounding Available''
state and suppresses the LLM output entirely, preventing operators from acting
on fabricated industrial context.

% ════════════════════════════════════════════════════════════════════════════
\section{Performance, Evaluation, and Trade-offs}
% ════════════════════════════════════════════════════════════════════════════

\begin{table}[htbp]
\centering
\caption{OmniWatch Performance Metrics Summary}
\label{tab:perf}
\begin{tabular}{@{}lll@{}}
\toprule
\textbf{Metric} & \textbf{Value} & \textbf{Notes} \\
\midrule
End-to-end triage latency  & 15--32\,s & Operator wait time (STARK + LLM) \\
STARK proof generation     & 9--17\,s  & RISC Zero local CPU prover \\
Phi-3 Mini RAG throughput  & 10--25 tokens/s & INT4 quantised, CPU \\
Raw STARK receipt size     & $\sim$217--250\,KB & Bincode $\rightarrow$ base64 \\
Ingestion memory footprint & 50--200\,MB & Polars batched reader \\
Technology Readiness Level & TRL 4--5  & Laboratory + simulated OT testbed \\
\bottomrule
\end{tabular}
\end{table}

The measured 15--32 second triage latency validates the project's core
research hypothesis: current verifiable computation and local generative
AI are too slow for real-time automated mitigation. This empirically justifies
the necessity of cryptographically signed human-in-the-loop authorization
as the superior operational control for post-detection governance.

% ════════════════════════════════════════════════════════════════════════════
\section{Commercial and Regulatory Alignment}
% ════════════════════════════════════════════════════════════════════════════

The platform explicitly aligns with the logging, monitoring, and
audit-trail objectives of NCA OTCC-1:2022 Subdomain~2-11
(Cybersecurity Event Logs and Monitoring Management). By requiring
mathematically verified human authorization before any automated
remediation action, OmniWatch addresses the stringent operational safety
requirements of critical water utility infrastructure and Vision~2030
smart city deployments in the Kingdom of Saudi Arabia.

% ════════════════════════════════════════════════════════════════════════════
\section{Conclusion}
% ════════════════════════════════════════════════════════════════════════════

OmniWatch demonstrates that zero-knowledge proof systems can be
practically integrated into OT/ICS security operations. The Split-Brain
architecture resolves the fundamental tension between floating-point
machine learning and the integer-only constraints of the RV32IM ISA.
By combining a dual-verification SOAR gate with a triple-layer audit
trail, the platform offers a robust, commercially aligned deployment
posture that effectively transitions threat response from probabilistic
black-box automation into mathematically deterministic, human-supervised
governance.

% ── Bibliography placeholder ───────────────────────────────────────────────
\newpage
\section*{References}
\addcontentsline{toc}{section}{References}

\begin{enumerate}
  \item RISC Zero, Inc., ``RISC Zero zkVM Documentation,''
        \url{https://dev.risczero.com/}, 2024.
  \item Canadian Institute for Cybersecurity, ``CIC-IDS-2017 Dataset,''
        University of New Brunswick, 2017.
  \item Splunk, Inc., ``Boss of the SOC Dataset Version 3 (BOTSv3),'' 2019.
  \item National Cybersecurity Authority (NCA), Kingdom of Saudi Arabia,
        ``Operational Technology Cybersecurity Controls (OTCC-1:2022),'' 2022.
  \item C.~Pevny, ``Loda: Lightweight on-line detector of anomalies,''
        \textit{Machine Learning}, vol.~102, no.~2, pp.~275--304, 2016.
  \item E.~S.~Page, ``Continuous inspection schemes,''
        \textit{Biometrika}, vol.~41, no.~1/2, pp.~100--115, 1954.
  \item C.~Masson, J.~E.~Rim, and H.~K.~Lee, ``DDSketch: A fast and
        fully-mergeable quantile sketch with relative-error guarantees,''
        \textit{Proc.\ VLDB Endowment}, vol.~12, no.~12, 2019.
  \item W3C and FIDO Alliance, ``Web Authentication: An API for accessing
        Public Key Credentials (Level 2),'' 2021.
\end{enumerate}

\end{document}
